% SHARD-ID: CA-41-QUBIT-2D-POINCARE-EQUATION-SCHEMA
% SHARD-TITLE: Two-Dimensional Qubit Poincare Equation Schema
% SHARD-SUMMARY: Gives a proposal-level square-lattice edge-density schema for ramp momenta, conservation, commuting momenta, boosts, and rotations.
% SHARD-SUMMARY: Every residual is a finite Pauli coefficient equation modulo a named two-dimensional lattice divergence.
% SHARD-SUMMARY: The shard records the orientation and midpoint choices needed before a two-dimensional checker can be built.
% SHARD-KEYWORDS: qubit, square lattice, 2+1 dimensions, Poincare, edge density, divergence

\section{Two-Dimensional Qubit Poincare Equation Schema}
\label{sec:qubit-2d-poincare-equation-schema}

\subsection{Status, Sources, and Dependencies}

\textbf{Status: proposal-level local schema.}  This is a finite algebraic
target for a future checker.  It is not yet a canonical stress tensor or a
continuum \(2+1\)-dimensional Poincare theorem.

Dependencies: CA-13, CA-36, CA-38, CONVENTIONS.md (o).

% Source: literature/md/2010.11121/2010.11121.md:204--222 -- local algebras
%   on hypercubic lattices and quasi-local algebra.
% Convention: CONVENTIONS.md (o) -- proposal-level square-lattice edge
%   orientation and divergence policy.

\subsection{Edge Densities and Vertex Currents}

On \(\mathbb Z^2\), write positively oriented edge densities
\[
  h^a_r=\sum_{\alpha,\beta=0}^3
  c^a_{\alpha\beta}\,
  \sigma_r^\alpha\sigma_{r+\hat a}^\beta,
  \qquad a\in\{x,y\}.
\]
At a vertex \(v\), order incident edges as \(W,S,E,N\), with midpoint offsets
\[
  m_W=(-1/2,0),\quad m_S=(0,-1/2),\quad
  m_E=(1/2,0),\quad m_N=(0,1/2).
\]
Use shared-leg coefficient tensors
\[
  q^W_{\lambda\mu}=c^x_{\lambda\mu},\quad
  q^E_{\lambda\mu}=c^x_{\mu\lambda},\quad
  q^S_{\lambda\mu}=c^y_{\lambda\mu},\quad
  q^N_{\lambda\mu}=c^y_{\mu\lambda}.
\]
For ordered incident pairs \(A<B\) in \(W<S<E<N\), the local commutator
coefficient on outer leg \(A\), center \(v\), outer leg \(B\) is
\[
  J^{AB}_{\lambda\kappa\rho}
  =
  -2\sum_{\mu,\nu=1}^3
  q^A_{\lambda\mu}q^B_{\rho\nu}\epsilon_{\mu\nu\kappa}.
  \tag{41.1}
\]
The ramp-generated momentum densities are
\[
  p_a(v)=\sum_{A<B}(m_B^a-m_A^a)J^{AB}(v).
  \tag{41.2}
\]
For \(AB=WS,WE,WN,SE,SN,EN\), the \(x\)-weights are
\[
  1/2,\ 1,\ 1/2,\ 1/2,\ 0,\ -1/2,
\]
and the \(y\)-weights are
\[
  -1/2,\ 0,\ 1/2,\ 1/2,\ 1,\ 1/2.
\]

\subsection{Divergence-Relaxed Residuals}

Every residual below is a finite Pauli tensor \(R\) on a named patch.  The
strong filter is \(R=0\).  The default bulk filter is the existence of
finite-support tensors \(U_x,U_y\) with
\[
  R=(1-\tau_{-\hat x})U_x+(1-\tau_{-\hat y})U_y.
  \tag{41.3}
\]
Coefficientwise, this is an explicit linear equation in the auxiliary
divergence tensors and a polynomial equation in \(c^x,c^y\), using the Pauli
product constants of CA-38.

\subsection{Necessary Bulk Equations}

Momentum conservation is
\[
  C_a:=
  \sum_{e:\operatorname{supp}h_e\cap\operatorname{supp}p_a(0)\neq\varnothing}
  i[h_e,p_a(0)]
  =\nabla\cdot U^{(C_a)},
  \qquad a=x,y.
  \tag{41.4}
\]
The commuting-translation equation is
\[
  T_{xy}:=
  \sum_{r:\operatorname{supp}p_x(0)\cap\operatorname{supp}p_y(r)\neq\varnothing}
  i[p_x(0),p_y(r)]
  =\nabla\cdot U^{(T)}.
  \tag{41.5}
\]
Let \(e_0=h^x_0+h^y_0\) denote the chosen local energy cell.  With
\(m_e^a\) the \(a\)-coordinate of the edge midpoint,
\[
  B_{ab}:=
  \sum_{e:\operatorname{supp}h_e\cap\operatorname{supp}p_b(0)\neq\varnothing}
  m_e^a\,i[p_b(0),h_e]
  -\delta_{ab}v_a^2 e_0
  =\nabla\cdot U^{(B_{ab})}.
  \tag{41.6}
\]
This is the \(i[P_b,K_a]=\delta_{ab}v_a^2H\) equation after the coordinate
weighted conservation term is removed by (41.4).

If one defines the rotation candidate
\[
  J=\sum_v(v_xp_y(v)-v_yp_x(v)),
\]
then the translation-rotation residual is
\[
  R_c=
  \sum_r\left(r_x\,i[p_c(0),p_y(r)]
        -r_y\,i[p_c(0),p_x(r)]\right)
  -(\delta_{cx}p_y(0)-\delta_{cy}p_x(0))
  =\nabla\cdot U^{(R_c)}.
  \tag{41.7}
\]
This last equation is quartic in \(c^x,c^y\) before the quadratic subtraction.
It depends on the chosen momentum-density split, so it is less canonical than
the boost equations.

\subsection{Implementation Boundary}

CA-41 supplies the finite equations a checker should implement.  The next
implementation step is to enumerate the patch supports in (41.4)--(41.7),
embed every edge tensor into the common patch, and emit Pauli coefficients plus
the divergence auxiliaries.  Until that exists, the \(2+1\)-dimensional claims
remain proposal-level.

